\paragraph{Raw, projected, and tail summaries.}
We compare four summaries of retained validation geometry: mean raw task second moment $A$, arithmetic mean projected information $P$, reciprocal mean inverse information $H$, and inverse-CVaR certificate $C$. For the four equally weighted groups, let $F_g=\widehat F_g^{\rid}$ and $L_g=(F_g+\varepsilon)^{-1}$. Then
\[
\begin{aligned}
A&=\tfrac14\sum_g\widehat J_{TT,g}, &
P&=\tfrac14\sum_g F_g,\\
H&=\left(\tfrac14\sum_g L_g\right)^{-1}, &
C&=\left[\CVaR_{0.95,\pi}(L_g)\right]^{-1},\qquad \pi_g=\tfrac14.
\end{aligned}
\]
Thus $H$ is stabilized harmonic projected information, not mean inverse information. All four summaries use uncentered feature moments, nuisance ridge $10^{-3}$, and reciprocal stabilization $10^{-8}$. This post hoc analysis joins validation summaries to held-out interventions in the primary (seeds 22--26) and development (17--21) cohorts. It compares summaries of existing models, not training objectives.

\begin{table}[htbp]
\centering\small
\setlength{\tabcolsep}{3pt}
\caption{Four validation summaries and their intervention associations. Panels (a)--(b): means $\pm$ nominal 95\% normal-approximation half-widths, $1.96s/\sqrt5$; changes are paired within seed and drops are signed percentage points. Selected FG includes accuracy-budget fallbacks (primary seeds 22, 24; development seed 20). Panel (c): descriptive Spearman coefficients for ten models arranged in five ERM/FG seed pairs per cohort, with average ranks for tied outcomes.}
\label{tab:four-way-swap}
\textit{(a) Model summaries and paired changes}\par\smallskip
\begin{tabular}{@{}llcccc@{}}
\toprule
Cohort & Model/change & $A$: raw & $P$: projected & $H$: harmonic & $C$: inverse-CVaR\\
\midrule
Primary & ERM & $1.142\pm0.012$ & $0.500\pm0.010$ & $0.442\pm0.011$ & $0.289\pm0.014$\\
& Selected FG & $0.730\pm0.086$ & $0.484\pm0.047$ & $0.435\pm0.043$ & $0.330\pm0.043$\\
& FG $-$ ERM & $-0.412\pm0.079$ & $-0.016\pm0.046$ & $-0.007\pm0.041$ & $+0.041\pm0.046$\\
Development & ERM & $1.152\pm0.005$ & $0.489\pm0.011$ & $0.413\pm0.014$ & $0.255\pm0.018$\\
& Selected FG & $0.754\pm0.118$ & $0.519\pm0.065$ & $0.433\pm0.057$ & $0.303\pm0.082$\\
& FG $-$ ERM & $-0.398\pm0.118$ & $+0.030\pm0.063$ & $+0.020\pm0.066$ & $+0.048\pm0.095$\\
\bottomrule
\end{tabular}
\par\smallskip
\textit{(b) ERM group-selection rules and controlled drops}\par\smallskip
\begin{tabular}{@{}>{\raggedright\arraybackslash}p{.34\linewidth}llrr@{}}
\toprule
Rule & Candidates & Selected, each cohort & Primary & Development\\
\midrule
Minimum raw second moment & All four & LL 4/5; LW 1/5 & $+15.8\pm21.3$ & $+16.9\pm23.2$\\
Minimum projected / largest inverse contribution / CVaR-active group & All four & LW 5/5 & $-26.9\pm2.0$ & $-29.7\pm2.2$\\
Largest inverse contribution & Eligible & LL 5/5 & $+26.7\pm1.8$ & $+29.2\pm2.8$\\
\bottomrule
\end{tabular}
\par\smallskip
\textit{(c) Model-level rank associations with held-out outcomes}\par\smallskip
\begin{tabular}{@{}lrrrr@{}}
\toprule
& \multicolumn{2}{c}{Primary} & \multicolumn{2}{c}{Development}\\
Summary & $D_{\max}$ & $W_{\mathrm{opp}}$ & $D_{\max}$ & $W_{\mathrm{opp}}$\\
\midrule
$A$: mean raw & $-0.10$ & $+0.26$ & $+0.26$ & $-0.13$\\
$P$: mean projected & $-0.09$ & $+0.21$ & $-0.03$ & $-0.43$\\
$H$: reciprocal mean inverse & $-0.15$ & $+0.21$ & $-0.12$ & $-0.40$\\
$C$: inverse-CVaR & $+0.15$ & $-0.24$ & $-0.03$ & $-0.48$\\
\bottomrule
\end{tabular}
\end{table}

For Table~\ref{tab:four-way-swap}(c), $D_{\max}=100\max_g(\operatorname{Acc}_{\mathrm{same},g}-\operatorname{Acc}_{\mathrm{opposite},g})$ and $W_{\mathrm{opp}}=\min_g\operatorname{Acc}_{\mathrm{opposite},g}$, both over all four groups. Integer correct counts recover exact outcome ties: primary ERM seeds 22 and 25 share drop $348/642$, despite different floating-point subtractions. LL and LW denote landbird/land and landbird/water. Negative drops mean opposite-type recomposition improves accuracy. The largest group contributes $38.3\%$ and $40.6\%$ of summed inverse information on average in the primary and development cohorts, respectively. The original headline rules differed: surviving fraction in the primary cohort, absolute projected information in development; this comparison consistently uses all four groups for model summaries and correlations.

Projection and aggregation give different model summaries, with rank associations varying across cohorts. For four equally weighted groups and $\alpha=0.95$, $C=\min_gF_g+\varepsilon$; its active group also minimizes projected information and maximizes inverse contribution. Here that group is LW, whose opposite-type recomposition improves accuracy. The prospective WW diagnosis instead uses the surviving fraction within the eligible set. This comparison does not establish superior swap-damage prediction by inverse-CVaR; paired changes in $P,H,C$ all have intervals containing zero. The artifact provides per-model and per-seed values, exact-count outcomes, input hashes, and replay code.
